\documentclass[12pt]{article}
\usepackage[utf8]{inputenc}
\usepackage[T1]{fontenc}
\usepackage{hyperref}
\title{Ubykh Language Resources for K-12 Curriculum Development}
\author{Compiled by Claude Sonnet 5 (claude-sonnet-5) for Andr\'as Kornai}
\date{}
\begin{document}
\maketitle

\section{Introduction}
Ubykh (ISO 639-3 uby; Glottocode ubyk1235) is an extinct Northwest Caucasian (Abkhaz-Adyge) language, formerly spoken in Turkey. The last fluent speaker, Tevfik Esen\c{c}, died on 7 October 1992. Linguists including Dum\'ezil, Vogt and Hewitt recorded him and made extensive notes. Any curriculum would be heritage and revitalization work built on documentation, not on living speakers.

\section{Primary Educational Resources}
No school primers or textbooks were found in the searches. The Colarusso sampler below is an annotated collection of texts and could serve as source material for teaching.
\begin{itemize}
\item John Colarusso, \emph{An Ubykh Sampler} (compiled, translated and analyzed): \url{https://aheku.net/f1jok/files/4054/1-anubykhsampler.pdf}
\end{itemize}

\section{Linguistic and Reference Materials}
\begin{itemize}
\item Wikipedia overview of Ubykh grammar (polysynthetic, agglutinative, ergative-absolutive): \url{https://en.wikipedia.org/wiki/Ubykh_grammar}
\item Hans Vogt's 1963 Ubykh dictionary, later revised by Dum\'ezil and Charachidz\'e, is described in a secondary source. The dictionary itself was not seen. \url{https://progressivegeographies.com/2023/04/28/indo-european-thought-in-twentieth-century-france-update-11-dumezil-and-charachidzes-work-on-ubykh-levi-strauss-and-his-archive-eliades-correspondence-koyr}
\item Dum\'ezil and Charachidz\'e's \emph{Le Verbe Oubykh} (1975) is mentioned in the same source. The book itself was not seen.
\item Charachidz\'e's planned new dictionary remained unfinished and unpublished at his death in 2010: \url{https://abkhazworld.com/aw/abkhazians/language/567-in-memoriam-georges-charachidze}
\item An English grammar by Rhona Fenwick is mentioned in a secondary source. No link was found.
\end{itemize}

\section{Religious and Community Resources}
No community organisation or religious materials specific to Ubykh were found in the searches.

\section{Digital and Audio Resources}
\begin{itemize}
\item Omniglot page on Ubykh (alphabet and background): \url{https://www.omniglot.com/writing/ubykh.htm}
\end{itemize}
Recordings by Dum\'ezil, Vogt and Hewitt are said to exist, but no archive location was found in the searches.

\section{Curriculum Development Context}
There are no native speakers, so any curriculum is a heritage and revitalization programme built from archival recordings, texts and grammars. The first steps are to locate the recordings and the Vogt and Dum\'ezil publications through library catalogues. The community connection is with Ubykh-descended groups in Turkey. No specific contacts were found.
\end{document}